El proyecto nació de la idea del Ing. Wilson Iván Félix Reyna, primer jefe del autor, de construir una librería que combinara inteligencia artificial y ciberseguridad. Se desarrolló entre mayo y octubre de 2026 junto con Sebastián Huertas Gómez, quien implementó el extractor inicial y el lado npm de la librería y desarrolló el componente de ciberseguridad, mientras que el autor desarrolló el componente de aprendizaje automático. Por recomendación del asesor, la evaluación se amplió a cuatro aplicaciones de terceros. Durante la redacción, la verificación de cada cifra contra el repositorio llevó a corregir afirmaciones previas, y esas correcciones y sus fuentes se documentan en los anexos.

Este documento reporta el componente de aprendizaje automático de logSguarDian, desarrollado junto con Sebastián Huertas Gómez, cuya tesis aborda el componente de ciberseguridad y la librería. El trabajo se limita a cuatro categorías de ataque intra-petición, a un entorno de prueba controlado y a aplicaciones de terceros evaluadas con muestras pequeñas. No incluye la validación con una startup guatemalteca prevista en el protocolo, que era opcional. Los modelos cuyas métricas se reportan (\texttt{rf\_v11} e \texttt{if\_v10}) no están publicados. Las limitaciones que acotan las cifras principales se declaran en los capítulos de Resultados y Discusión y en los anexos.

\textbf{Declaración de uso de inteligencia artificial.} El autor concibió el proyecto y su enfoque, seleccionó los modelos, concibió el pipeline de MLOps y el análisis de interpretabilidad con SHAP, y definió el alcance y los resultados que se reportan. Para el código de entrenamiento y los scripts de análisis, la redacción y la verificación de coherencia entre capítulos, trabajó con Claude, de Anthropic, mediante conversación y mediante un agente de código. El autor dirigió ese trabajo y asume la responsabilidad del contenido.

El autor agradece al Ing. Wilson Iván Félix Reyna, su primer jefe en su primer trabajo, quien propuso la idea original de una librería que combine inteligencia artificial y ciberseguridad. Agradece al Ing. Erick Francisco Marroquín Rodríguez, asesor de este trabajo, por su guía y por la recomendación de evaluar el sistema en aplicaciones de terceros. Agradece al Ing. Bacilio Alexander Bolaños Lima, catedrático, quien lo acompañó durante el desarrollo de este proyecto, y a Sebastián Huertas Gómez, compañero de proyecto.

Agradece también a los catedráticos que lo formaron durante la carrera y a la Universidad del Valle de Guatemala, su casa de estudios. Agradece a sus compañeros de carrera y a sus amigos por el acompañamiento durante estos años. Agradece a sus padres y a su hermana por su apoyo. Agradece a las comunidades de código abierto cuyos conjuntos de datos y aplicaciones vulnerables hicieron posible la evaluación, entre ellos SecLists, PayloadsAllTheThings, OWASP Juice Shop, DVNA, dvws-node y node-api-goat.
